This paper proposes that language has a limit, and that the limit is generative. Across six dimensions of relational life, its existence, its maintenance, the value it generates, its ethics, its justice, and its evolution, there is in each case a core that language reaches toward and cannot reach. The paper sets out these six limits, each as a claim of its own, and then argues that what stands at the limit is drive in the psychoanalytic sense. Drive is not a content language omits and might recover; it is the phenomenology of manque, the gap between the symbolic and the Real that resists complete symbolization. This gap has two faces, a pre-symbolic remainder cut off when the subject enters language and a resulting impossibility of closure, and they are one. Around it drive circles and desire slides. By a thought experiment the paper shows that closing the gap would extinguish rather than fulfill them, so that the non-closure of language is the condition of their persistence. It then follows drive into the time of relation and draws out the generativity of the limit in the poem, in an ethics without possession, and in a justice that guards what it cannot settle, with a coalgebraic and reduction-theoretic counterpart offered as a heuristic. The paper offers no final synthesis; its form enacts its content.
